\documentstyle[% 


righttag% 


,11pt% 


,amssymb% 


,russian%               remove in Eng. 


%,izvru%                 Set izven% in Eng. 


,izvru-ks%             for brief communications 


]{amsart} 


 


%  


 


\newcommand{\la}{\langle} 


\newcommand{\ra}{\rangle} 


\newcommand{\myza}{\Bbb Z[\alpha]} 


\newcommand{\myzx}{\Bbb Z\langle x\rangle} 


\newcommand{\Mod}{\operatorname{mod}} 


\newcommand{\tts}[1]{} 


 


%\hoffset=-15mm%                  For Eng. 


 


\setcounter{page}{81} 


\god{1999} 


\nomer{11~(450)} %                        Remove in Eng. 


%\nomer{Vol.\,43, No.\,11}%               Set in Eng. 


%\str{\pageref{begin}--\pageref{end}}%  Set in Eng. 


\udk{512.552.7} 


\author{\it Р.Ж.\,Алеев, Г.А.\,Панина} 


% NB:   ^^ remove in Eng. 


\title{Единицы циклических групп порядков 7 и 9} 


 


\date{} 


% \pagestyle{myheadings}%         Set in Eng. 


\pagestyle{plain} %              Remove in Eng. 


\begin{document} 


% \thispagestyle{plain}%          Set in Eng. 


\maketitle 


\label{begin} 


 


\section{Введение} 


 


Данная работа написана на основе доклада авторов на Международной 


конференции ``Алгебра и анализ'', посвященной 100-летию со дня 


рождения  


Н.Г.\,Чеботарева (Казань, июнь 1994 г.) [1]. 


В работе будут рассматриваться целочисленные групповые кольца,  


и для краткости будем называть их единицы ($=$\,обратимые элементы)  


единицами соответствующих групп. 


Полная информация о группах единиц циклических групп порядков 5,  


8, 10 и 12 приведена в [2]. В данной 


работе описаны группы единиц циклических групп  


порядков 7 и 9. Таким образом, получается описание групп единиц всех  


циклических групп порядка не более 10, потому что для других порядков  


единицы тривиальны. 


 


\section{Единицы циклической группы порядка 7} 


 


Пусть $\alpha$ --- первообразный корень степени 7 из единицы, 


$\eta_1=\alpha+\alpha^{-1}$ и 


$\eta_2=\alpha^2+\alpha^{-2}$. Пусть $\myza$ --- кольцо круговых целых 


степени 7, и $U(\myza)$ 


--- его группа единиц. 


 


\begin{lem}[{\series{m}\shape{n}\selectfont[3], c.\,241}] 


$U(\myza)=\la-1\ra\times\la\alpha\ra\times 


\la\eta_1\ra\times\la\eta_2\ra$. 


\end{lem} 


 


Пусть $\myzx$ --- целочисленное групповое кольцо циклической группы  


$\la x\ra$ порядка 7, и $U(\myzx)$ --- его группа единиц. Пусть 


\begin{alignat*}{2}


\varphi&:\myzx\to\myza &\ \; \varphi\bigg(\sum^6_{i=0} 


a_i x^i\bigg)&=\sum^6_{i=0}a_i\alpha^i, \\ 


% 


\psi&:\myzx\to\Bbb Z &\ \; \psi\bigg(\sum^6_{i=0} 


a_i x^i\bigg)&=\sum^6_{i=0}a_i. 


\end{alignat*} 


Эти отображения --- гомоморфизмы колец. 


 


\begin{lem} 


Если $u_1$ --- единица кольца $\myzx$ и $\varphi(u_1)=\eta^k_1$, то 


$k$ 


делится на $3$. Причем, если 


$\varphi(u_1)=\eta^3_1$, то $u_1=-1+2x-x^2-x^5+2x^6$ 


и $u^{-1}_1=-3+x+3x^2-2x^3-2x^4+3x^5+x^6$. 


\end{lem} 


 


\begin{pf} 


Пусть $\varphi(u_1)=\eta_1$. Тогда $u_1\in x+x^6+\ker\varphi$. 


Так как $\ker\varphi=\Big\{\sum\limits^6_{i=0}a_i x^i\big|a_i=a_6$ для 


всех $i=1,\dotsc,5\Big\}$, то 


$u_1=x+x^6+a\sum\limits^6_{i=0}x^i$, 


но тогда $\psi(u_1)=2+7a=\pm1$, что невозможно. 


Пусть 


$\varphi(u_1)=\eta^3_1=\alpha^3+3\alpha+3\alpha^{-1}+\alpha^{-3}$. 


Тогда $u_1=3x+x^3+x^4+3x^6+a\sum\limits^6_{i=0}x^i$ и 


$\psi(u_1)=8+7a=\pm1$. 


Отсюда $a=-1$ и $u_1=3x+x^3+x^4+3x^6- 


\sum\limits^6_{i=0}x^i=-1+2x-x^2-x^5+2x^6$. 


\end{pf} 


 


Аналогично доказываются 


 


\begin{lem} 


Если $u_2$ --- единица кольца $\myzx$ и $\varphi(u_2)=\eta^n_2$, то 


$n$ делится на $3$. Причем если $\varphi(u_2)=\eta^3_2$, то 


$u_2=-1+2x^2-x^3-x^4+2x^5$ 


и $u^{-1}_2=-3-2x+x^2+3x^3+3x^4+x^5-2x^6$. 


\end{lem} 


 


\begin{lem} 


Если $u_3$ --- единица кольца $\myzx$, 


$\varphi(u_3)=\eta^r_1\eta^s_1$, $r$ и $s$ не 


делятся на $3$, то\linebreak $r\not\equiv s\,(\Mod 3)$. Причем если 


$\varphi(u_3)=\eta_1\eta^2_2$, то  


$u_3=-1+x+x^6$ и $u^{-1}_3=1+x-x^3-x^4+x^6$. 


\end{lem} 


 


\begin{thm} 


Пусть $U(\myzx)$ --- группа единиц целочисленного группового 


кольца циклической группы порядка $7$ и $\varphi$ определено как 


ранее.  


Тогда 


$$ 


U(\myzx)=\la-1\ra\times\la x\ra\times\la u_2\ra 


\times\la u_3\ra, 


$$ 


где $u_2$ и $u_3$, как в леммах $3$ и $4$. Кроме того, 


$U(\myza)/\varphi(U(\myzx))\cong\Bbb Z_3$. 


\end{thm} 


 


\begin{pf} 


Пусть $u$ --- произвольная единица кольца $\myzx$. 


Тогда по лемме 1\; $\varphi(u)=(-1)^i\alpha^j\eta^k_1\eta^n_2$ и 


$\varphi((-1)^i x^{-j}u)=\eta^k_1\eta^n_2$. Пусть  


$k=3p+r$ и $n=3t+s$, где $r,s\in\{0,1,2\}$, и по леммам 2--4 либо  


$r=s=0$, либо $rs=2$. По леммам 2 и 3\; $\varphi((-1)^i 


x^{-j}u^{-p}_1u^{-t}_2u)=\eta^r_1\eta^s_2$ и 


по лемме 4\; $\varphi((-1)^i x^{-j}u^{-p}_1u^{-t}_2u^{-1}_3u)=1$ при 


$r=1$ и $s=2$, а $\varphi((-1)^i 


x^{-j}u^{-p}_1u^{-t}_2u^{-1}_1u^{-1}_2u_3u)=1$ при $r=2$ и $s=1$. 


Так как по теореме 12 из [2]  


гомоморфизм $\varphi$ инъективен на группе единиц кольца $\myzx$, то 


$$ 


U(\myzx)=\la-1\ra\times\la x\ra\times\la u_1,u_2,u_3\ra. 


$$ 


Заметим, что $\eta^3_1=(\eta_1\eta^2_2)^3(\eta^3_2)^{-2}$ и в силу 


инъективности $\varphi$ на группе  


единиц $u_1=u^3_3u^{-2}_2$. Отсюда 


$\la u_1,u_2,u_3\ra=\la u_2,u_3\ra$ и 


\begin{multline*} 


U(\myza)/\varphi(U(\myzx))=(\la-1\ra\times\la\alpha\ra 


\times\la\eta_1\ra\times\la\eta_2\ra)/ 


(\la-1\ra\times\la\alpha\ra\times\la\eta^3_2\ra 


\times\la\eta_1\eta^2_2\ra)\cong \\ 


% 


\cong(\la\eta_1\ra\times\la\eta_2\ra)/ 


(\la\eta^3_2\ra\times\la\eta_1\eta^2_2\ra). 


\end{multline*} 


Откуда $U(\myza)/\varphi(U(\myzx))\cong\Bbb Z_3$. 


\end{pf} 


 


\section{Единицы циклической группы порядка 9} 


 


3.1. {\it Единицы круговых целых степени} 9. 


Пусть $\alpha$ --- первообразный корень степени 9 из единицы и 


$k=\Bbb Q(\alpha)$. 


Пусть $P=\la1-\alpha^i\mid i=1,\dotsc,8\ra$ --- подгруппа 


мультипликативной группы 


поля $k$. По определению подгруппой круговых единиц назовем подгруппу 


$C=P\cap U(\myza)$, где $U(\myza)$ --- группа единиц кольца круговых  


целых степени 9. 


 


\begin{lem} 


$U(\myza)=C$, т.\,е. всякая единица кольца $\myza$ является круговой 


единицей. 


\end{lem} 


 


{\bf Доказательство}\,\footnote{\,Первоначальное доказательство 


следовало доказательству леммы 1. Авторы  


признательны Л.В.\,Кузьмину на указание статьи [6].}. 


Пусть $U(\myza)=E$, $E^+$ --- подгруппа вещественных 


единиц кольца $\myza$ и $C^+=C\cap E^+$. Из ([6], сс.\,107, 119) 


следует,  


что $E^+/C^+\cong E/C$ и $|E^+/C^+|=h^+$, где $h^+$ --- число классов 


максимального 


вещественного подполя поля $k$. Согласно [7]\; $h^+=1$ и потому  


$U(\myza)=E=C.\quad\square$


 


\begin{thm} 


Группа единиц кольца круговых целых степени $9$ 


$$ 


U(\myza)=\la-1\ra\times\la\alpha\ra\times\la\lambda\ra 


\times\la\mu\ra, 


$$ 


где $\lambda=\alpha+\alpha^8$ и $\mu=\alpha^4+\alpha^5$. 


\end{thm} 


 


\begin{pf} 


Для любого $j$\; $1-\alpha^{9-j}=(-\alpha^{-j})(1-\alpha^j)$ и 


$(1-\alpha)(1-\alpha^2)(1-\alpha^4)=(-\alpha^2)(1-\alpha^3)$. Поэтому 


для $u\in C=U(\myza)$ 


$$ 


u=(-\alpha)^j(1-\alpha)^n(1-\alpha^2)^p(1-\alpha^4)^r 


$$ 


для подходящих целых $j$, $n$, $p$, $r$. 


 


Пусть $N$ --- норма расширения $k/\Bbb Q$. По лемме 1.3 из ([4], 


с.\,19) 


$N(1-\alpha^j)=3$ для $(j,3)=1$. Так как $u$ --- единица, то 


$N(u)=1^j\cdot3^n\cdot3^p\cdot3^r=3^{n+p+r}=1$ и $n+p+r=0$. 


Поэтому 


$$ 


u=(-\alpha)^j\bigg(\frac{1-\alpha^2}{1-\alpha}\bigg)^p 


\bigg(\frac{1-\alpha^4}{1-\alpha}\bigg)^r. 


$$ 


Так как $\frac{1-\alpha^2}{1-\alpha}=1+\alpha$ и 


$\frac{1-\alpha^4}{1-\alpha}=(1+\alpha)(1+\alpha^2)$, а 


$\lambda=\alpha+\alpha^8=\alpha^{-1}(1+\alpha^2)$ и 


$\mu=\alpha^4+\alpha^5=\alpha^4(1+\alpha)$, то 


$U(\myza)=\la-1\ra\times\la\alpha\ra\times\la\lambda\ra 


\times\la\mu\ra$. 


\end{pf} 


 


3.2. {\it Единицы кольца} $\myzx$. 


Для дальнейших вычислений удобно ввести следующие отображения: 


$$ 


\varphi:\myzx\to\myza,\quad 


\tau:\myzx\to\Bbb Z[\alpha^3],\quad 


\psi:\myzx\to\Bbb Z, 


$$ 


действующие по правилам 


$$ 


\varphi\bigg(\sum^8_{i=0}a_i x^i\bigg)= 


\sum^8_{i=0}a_i\alpha^i,\quad 


\tau\bigg(\sum^8_{i=0}a_i x^i\bigg)= 


\sum^8_{i=0}a_i\alpha^{3i},\quad 


\psi\bigg(\sum^8_{i=0}a_i x^i\bigg)= 


\sum^8_{i=0}a_i. 


$$ 


Эти отображения --- гомоморфизмы колец. 


\medskip 


 


{\bf Лемма 6.} 


 


1) $\ker\varphi=\Big\{\sum\limits^2_{i=0}c_i x^i(1+x^3+x^6)\mid 


c_0,c_1,c_2\in\Bbb Z\Big\}$; 


 


2) $U(\Bbb Z[\alpha^3])=\{\pm1,\pm\alpha^3,\pm(1+\alpha^3)\}$. 


 


\begin{pf} 


1) Стандартно.  


 


2) Известно (см., напр., [5], гл.\,II, \S\,7). 


\end{pf} 


\addtocounter{lem}{1} 


 


\begin{lem} 


Если $u_1$ --- единица кольца $\myzx$ и $\varphi(u_1)=\lambda^k$, то 


$k$ 


делится на $3$. Причем если $\varphi(u_1)=\lambda^3$, то 


$u_1=-1+2(x+x^8)-(x^2+x^7)-(x^4+x^5)$ и 


$u^{-1}_1=5+(x+x^8)-5(x^2+x^7)-3(x^3+x^6)+4(x^4+x^5)$. 


\end{lem} 


 


{\bf Доказательство.} Пусть $\varphi(u_1)=\lambda$. Тогда 


$u_1=x+x^8+c_0(1+x^3+x^6)+c_1(x+x^4+x^7)+c_2(x^2+x^5+x^8)$, 


но отсюда $\psi(u_1)=2+3(c_0+c_1+c_3)=\pm1$ 


и потому $c_0+c_1+c_2=-1$. 


Легко видеть, что $\tau(u_1)=(3c_0-3c_2-1)+\alpha^3(3c_1-3c_2)$. 


По предыдущей лемме $c_1=c_2$, $3c_1-3c_2-1=-1$ и получается 


противоречие с $c_0+c_1+c_2=-1$. Поэтому нет единицы $u_1$ с 


$\varphi(u_1)=\lambda$. Если 


$\varphi(u_1)=\lambda^3=\alpha^3+3\alpha+3\alpha^{-1}+\alpha^{-3}$, то 


$u_1=3x+x^3+x^6+3x^8+c_0(1+x^3+x^6)+c_1(x+x^4+x^7)+c_2(x^2+x^5+x^8)$ и 


$\psi(u_1)=8+3(c_0+c_1+c_2)=\pm1$. 


Отсюда $c_0+c_1+c_2=-3$. Теперь 


$\tau(u_1)=(3c_0-3c_2-1)+\alpha^3(3c_1-3c_2)$. 


По предыдущей лемме $c_1=c_2$, $3c_1-3c_2-1=-1$ и $c_0=c_1=c_2=-1$. 


Таким образом, получаем элемент 


\begin{multline*} 


u_1=3x+x^3+x^6+3x^8-(1+x^3+x^6)-(x+x^4+x^7)- 


(x^2+x^5+x^8) = \\ 


% 


=-1+2(x+x^8)-(x^2+x^7)-(x^4+x^5).\quad\square 


\end{multline*} 


 


Аналогично доказываются 


 


\begin{lem} 


Если $u_2$ --- единица кольца $\myzx$ и $\varphi(u_2)=\mu^n$, то $n$ 


делится на $3$. Причем если $\varphi(u_2)=\mu^3$, то 


$u_2=-1-(x+x^8)-(x^2+x^7)+2(x^4+x^5)$ и 


$u^{-1}_2=5-5(x+x^8)+4(x^2+x^7)-3(x^3+x^6)+(x^4+x^5)$. 


\end{lem} 


 


\begin{lem} 


Если $u_3$ --- единица кольца $\myzx$, $\varphi(u_3)=\lambda^r\mu^s$, 


$r$ и $s$ 


не делятся на $3$, то\linebreak $r\not\equiv s\,(\Mod 3)$. Причем, если 


$\varphi(u_3)=\lambda\mu^2$, то 


$u_3=1+(x+x^8)-(x^3+x^6)-(x^4+x^5)$ и 


$u^{-1}_3=-1+(x+x^8)-(x^2+x^7)$. 


\end{lem} 


 


Аналогично теореме 1 доказывается 


 


\begin{thm} 


Пусть $U(\myzx)$ --- группа единиц целочисленного группового 


кольца циклической группы порядка $9$ и $\varphi$ определено как 


ранее. Тогда 


$U(\myzx)=\la-1\ra\times\la x\ra\times\la u_2\ra\times\la u_3\ra$, 


где $u_2$ и $u_3$, как в леммах $8$ и $9$. Кроме того, 


$U(\myza)/\varphi(U(\myzx))\cong\Bbb Z_3$. 


\end{thm} 


 


\begin{thebibliography}{9} 


 


\bibitem{1} 


Алеев Р.Ж., Панина Г.А. {\em О единицах групповых колец циклических 


групп} // Алгебра и анализ. Тез. докл. международн. научн. конф., 


посв.  


100-летию со дня рожд. Н.Г.\,Чеботарева. Ч.\,I. -- Казань, 1994. -- 


С.\,6. 


\bibitem{2} 


Aleev R.\v Z. {\em Higman's central unit theory, units of integral 


group rings 


of finite cyclic groups and Fibonacci numbers} // Int. J. Algebra and 


Comput. -- 1994. -- V.\,4. -- \No\,3. -- P.\,309--358. 


\bibitem{3} 


Эдвардс Г. {\em Последняя теорема Ферма. Генетическое введение в 


алгебраическую теорию чисел}. -- М.: Мир, 1980. -- 484 с. 


\bibitem{4} 


Бовди А.А. {\em Мультипликативная группа целочисленного группового 


кольца}. -- Ужг. гос. ун-т. -- Ужгород, 1987. -- 210 c. -- 


Деп. в УкрНИИНТИ 24.09.87, \No\,2712-Ук87. 


\bibitem{5} 


Боревич З.И., Шафаревич И.Р. {\em Теория чисел}. -- 3-е изд. -- М.: 


Наука, 1985. -- 503 с. 


\bibitem{6} 


Sinnott W. {\em On the Stickelberger ideal and the circular units of a 


cyclomatic field} // Ann. Math. -- 1978. -- V.\,108. -- 


\No\,1. -- P.\,108--134. 


\bibitem{7} 


Masley J.M., Montgomery H.L. {\em Cyclotomic fields with unique 


factorization} // J. Reine und Angew. Math. -- 1976. -- Bd.\,286/287. 


-- S.\,248--256. 


 


\end{thebibliography} 


 


\vskip 2cm 


 


\postup{\it Челябинский государственный университет}% 


{$\qquad$\it Поступили} 


\smallskip


\postup{\it Медногорский электротехнический}% 


{{\it полный текст} 01.09.1994} 


\postup{\it завод АО ``Уралэлектро''}% 


{{\it краткое сообщение} 26.04.1999} 


\label{end} 


\end{document} 


